\scanpage{43}
\mspage{21}
prenant pour $\pi_{K'}$ le groupe unité (i.e. $K'=k$), on trouverait un hom. $K\longrightarrow k$ correspondant!

Il faut pour le moins, pour travailler à l'aise à partir d'homomorphismes $\pi_{K'}-\pi_K$ (au lieu de $E_{K'}-E_K$) supposer que le centralisateur dans $\pi_K$ de l'image de tt sous-groupe ouvert de $\pi_{K'}$ soit réduit à $(1)$ --- on dira que l'homomorphisme en question est \uline{anabélien} disons --- de telle façon qu'à partir de cet hom. (commutant à $\Gamma$) on reconstitue l'hom. d'extensions $E_K$ et $E_{K'}$, qui est l'objet vraiment essentiel. P. ex. si justement $K'=k$, donc $E_{K'}=\Gamma$, ce \unsure{que|qui} nous intéresse, ce ne sont pas les $\Gamma$-hom. de $\pi_{K'}=(1)$ (!) dans $\pi_K$, mais bien les \uline{sections} de $E_K$ sur $\Gamma$.

\uline{Question-conjecture} Soient $K,K'$ deux corps, extensions \inserted{de type} \unsure{finies|finie} de $\mathbb Q$, \deleted{considérons} \inserted{et} un morphisme $B_{K'}\longrightarrow B_K$ de topos sur $B_{\mathbb Q}$.
\margin{42}
